\section{Solved readout gates with a positive acceptance cost}\label{sec:solved}
Endpoint optimality alone says little about the boundary of an acceptance set. Here both that boundary and the remaining mode optimization can be solved. There is one cartridge remaining, two continuous positive terminal payoffs $G_1(R),G_2(R)$, a fixed current prior $\varpi$, and an acceptance multiplier $\alpha\in(0,1)$, with multiplier one on censoring. This is the earlier criterion with $\lambda=1$, $r_a(x)=\log\alpha<0$ and $r_a^\dagger=0$. Thus acceptance has a strictly positive cost, identical on every accepted raw component. Gates lie in $[0,1]$.

Write
\[
 U_j(x)=\int G_j(R)k_R^a(x)\varpi(dR),\quad
 M_{j,k}=\int R^kG_j(R)\varpi(dR)\quad(0\le k\le3),
\]
and let $\bar j=3-j$. Define
\begin{align*}
 S_j&=\pi M_{j,2}/a_0,\qquad
 N_j=M_{j,0}-S_j-(2T/a_0)M_{j,1},\\
 A_j&=(2T/a_0)M_{j,1},\qquad B_j=(2T\epsilon/a_0)M_{j,3}.
\end{align*}
Here $\pi$ denotes the circle constant and $\varpi$ denotes the current prior. On collision, $U_j(1,y)=A_j+B_j\cos(y-\theta)$; placement and no-hit components have values $S_j,N_j$.

For real $A,B$ set $J(A,B)=\int_0^{2\pi}(A+B\cos y)_+\,dy/(2\pi)$. Explicitly,
\begin{equation}
 J(A,B)=
 \begin{cases}
 A,&A\ge|B|,\\
 0,&A\le-|B|,\\
 \{A\arccos(-A/|B|)+\sqrt{B^2-A^2}\}/\pi,&|A|<|B|.
 \end{cases}
 \label{eq:J}
\end{equation}
At $A=B=0$ both first two values are zero.

\begin{theorem}[One-arc gates and a four-candidate mode solution]\label{thm:arc}
For a fixed mode the optimal value is $\max_{j=1,2}Q_j$, where
\begin{equation}
 Q_j=M_{j,0}+(\alpha S_{\bar j}-S_j)_+
 +(\alpha N_{\bar j}-N_j)_+
 +J(\alpha A_{\bar j}-A_j,\alpha B_{\bar j}-B_j).
 \label{eq:solved-value}
\end{equation}
For a maximizing $j$, censoring is followed by terminal decision $j$ and the optimal gate is
\begin{equation}
 g_j(x)=\one_{\{\alpha U_{\bar j}(x)>U_j(x)\}}.
 \label{eq:solved-gate}
\end{equation}
The collision acceptance set is empty, the whole circle, or a single circular arc (up to null ties). Its boundary is the explicitly computed threshold of $\cos(y-\theta)$. The other raw components are each wholly accepted or rejected.

Over the original compact mode family, an optimum has $\epsilon=E$, $T\in\{T_-,T_+\}$, and any phase with the gate rotated accordingly. Hence eight posterior payoff moments and four values in \eqref{eq:solved-value} solve the last-cartridge command problem; no continuous gate search remains.
\end{theorem}
\begin{proof}
For a gate $g$, Bayes optimization of the terminal decision gives
\[
 Q(g)=\alpha\int g\max(U_1,U_2)\,d\nu
       +\max_{j=1,2}\int(1-g)U_j\,d\nu.
\]
Interchange the finite maximum and the supremum over gates. For a fixed failed-branch decision $j$, the expression is
$M_{j,0}+\int g[\alpha\max(U_1,U_2)-U_j]\,d\nu$.
Since $U_j>0$ and $\alpha<1$, the contribution $(\alpha-1)U_j$ is strictly negative. The positive part of the bracket is therefore exactly $(\alpha U_{\bar j}-U_j)_+$. Pointwise optimization gives the stated gate and value. A maximizing failed-branch decision is consistent with its optimized gate: if another failed-branch decision improved it, the joint maximum just computed would not be maximal.

On the hit component the switching expression is an affine function of one cosine. Its positive set is a circle arc or a degenerate case, and integration over that arc gives \eqref{eq:J}. On the other components the switching function is constant. This proves the formula including all failure accounting.

For fixed $T$ and $j$, the only $\epsilon$-dependent term is $J(A,\epsilon B)$ with $A$ independent of $\epsilon$. It is convex in $\epsilon$, being an integral of positive parts of affine functions, and even by the shift $y\mapsto y+\pi$. Every even convex function is nondecreasing on $[0,\infty)$, so the maximum uses $E$. For fixed $\epsilon$, the placement term is constant in $T$, the no-hit term is the positive part of an affine function of $T$, and the collision term is linear in positive $T$ by homogeneity of $J$. Each $Q_j$ is convex in $T$; a maximum on an interval occurs at an endpoint. The phase only rotates a full-circle integral and its gate. Evaluating the two failed decisions at the two time endpoints therefore completes mode optimization.
\end{proof}

\subsection{A certified strict advantage over all report-blind gates}
The example compares gates, not a free equilibrium experiment with a counted one. Fix $T=1/20$, $\epsilon=1$ and $\theta=0$. Let
\[
 \varpi_* =\frac{u}{l+u}\delta_l+\frac{l}{l+u}\delta_u,
 \qquad \alpha=999/1000,
\]
\[
 G_1(R)=1-\frac{R-(l+u)/2}{u-l},\qquad G_2(R)=2-G_1(R).
\]
Both payoffs lie in $[1/2,3/2]$. In particular their values at $(l,u)$ are $(3/2,1/2)$ and $(1/2,3/2)$. The prior is chosen so that $M_{1,1}=M_{2,1}$, not so that the mark is uninformative: $M_{2,3}>M_{1,3}$.

\begin{proposition}[A positive, exact gate-performance certificate]\label{prop:example}
In this example, failed-branch decision 2 is optimal. Its gate rejects the placement-failure atom, accepts all no-hit reports, and accepts a collision precisely when
\[
 \cos y<c_*,\qquad -0.10631<c_*<-0.10630.
\]
The optimal value exceeds that of \emph{every report-blind constant gate} $g(x)\equiv p$, $0\le p\le1$, by more than $8\cdot10^{-4}$. This strict advantage persists, exceeding $7\cdot10^{-4}$, for the full-support prior $(1-10^{-5})\varpi_*+10^{-5}\operatorname{Unif}(I)$, using the same explicitly given gate and terminal rule against all report-blind gates.
\end{proposition}
\begin{proof}
All quantities are given by the preceding closed formulas. Put $A=A_1=A_2$ and retain $B_1,B_2$. Then the two collision switching pairs are
$((\alpha-1)A,\alpha B_2-B_1)$ and
$((\alpha-1)A,\alpha B_1-B_2)$.
The latter has negative second coordinate and threshold
$c_*=-(\alpha-1)A/(\alpha B_1-B_2)$. Direct substitution, with rational interval bounds described below, gives
\begin{align*}
 1.02689&<Q_1<1.02691,&1.02748&<Q_2<1.02749,\\
 1.02667&<V_{\rm accept}<1.02669,&1.01086&<V_{\rm reject}<1.01088.
\end{align*}
A sharper evaluation of the same expressions gives
$Q_2-V_{\rm accept}>8/10000$. The placement and no-hit switching signs follow from $\alpha S_2>S_1$, $S_2>S_1$, $\alpha N_1>N_2$ and $N_1>N_2$. In particular the failure-2 gate has the asserted component decisions. Always-accept value is
\[
 V_{\rm accept}=\alpha\{S_2+N_1+A+(B_2-B_1)/\pi\},
\]
and always-reject value is $M_{1,0}$. The command value is convex in a constant gate $p$, either from the explicit two-decision formula or Theorem~\ref{thm:endpoint}'s convexity argument. Thus every report-blind gate is bounded by the better of these two endpoints. This proves the strict comparison, not just comparison to one arbitrarily selected blind gate.

For transparency, the displayed strict inequalities have an exact rational certificate, not a floating quadrature premise. Use
$\pi=16\arctan(1/5)-4\arctan(1/239)$ and bracket each arctangent by its 48-term alternating series and the next term. For rational $x\ge0$, bracket $\sqrt{x}$ by
\[
 \frac{\lfloor\sqrt{10^{60}x}\rfloor}{10^{30}}
 \quad\hbox{and}\quad
 \frac{\lfloor\sqrt{10^{60}x}\rfloor+1}{10^{30}},
\]
using integer square roots after multiplying numerator and denominator. For the interior branch of $J$, use
$\arccos z=\pi/2-\arctan(z/\sqrt{1-z^2})$; every arctangent argument in this certificate lies in $[0,1/4]$. Interval addition, multiplication and division propagate these rational bounds. The accompanying \path{tests/rational_gate_certificate.py} prints the rational endpoints and checks the strict rational inequalities. The Machin identity follows from the tangent addition formula with both sides in the appropriate principal angle range; the alternating remainders and integer-root inequalities justify every enclosure.

Finally, the complete one-step payoff lies in $[\alpha/2,3/2]$, of width $D=3/2-\alpha/2$. The value of every fixed policy changes by at most $D\,\TV(\varpi,\varpi_*)$, as does the supremum over the common blind policy class. The explicit gate's advantage therefore falls by at most $2D\cdot10^{-5}$. The same rational certificate proves the remaining gap exceeds $7/10000$. The perturbed prior has full support; it is not necessary to recompute an optimum to obtain this lower bound.
\end{proof}

This is a strict selection advantage under a declared observation cost. It is not described as a temporal adaptive advantage at a one-step horizon, nor as a comparison against every unrestricted alternative apparatus. The four-candidate mode theorem and the full-support perturbation make the example a checkable control result inside the original positive laboratory rather than a finite-grid surrogate for it.
